\section{模型检验与稳健性分析}
\label{sec:validation}

前述结果来自多日、多个决策节点的连续优化。对于这类时序模型，仅有“求解器返回最优”并不足以说明方案可信：一方面，购电、储能和光伏必须在每个 $10$ min 时段满足能量守恒与设备边界；另一方面，计划还必须遵守真实的信息到达顺序，不能在决策时使用未来才会发生的数据。因此，本节从物理可行性、时间因果性、结算一致性和参数稳健性四个层面重新检验主结果。

\subsection{物理可行性与数值一致性}

首先把各问最终方案逐时重新代入母线平衡式和储能状态递推式，并同步检查储电量区间、充放电互斥以及问题三的费用分解。这样做的目的，是把“优化器给出的变量”重新还原为一条真实运行轨迹，检验其是否在每一个时段都能被物理实现。主要检验结果列于表~\ref{tab:validation}。

\begin{table}[htbp]
    \caption{模型物理可行性与数值一致性检验}
    \label{tab:validation}
    \centering
    \zihao{-5}
    \begin{tabular}{l l l l}
        \toprule
        检验项目 & 最大偏差或范围 & 判定标准 & 结论 \\
        \midrule
        母线能量平衡残差 & $<10^{-12}$ kWh & $10^{-7}$ kWh & 通过 \\
        储电量递推残差 & $<10^{-12}$ kWh & $10^{-7}$ kWh & 通过 \\
        储电量取值范围 & $[1200,\ 10800]$ kWh & 题设区间 & 通过 \\
        主方案同时充放电量 & $<10^{-8}$ kWh & $10^{-5}$ kWh & 通过 \\
        问题三费用回代差 & $<10^{-8}$ 元 & $10^{-7}$ 元 & 通过 \\
        结果工作簿回读差异 & $<10^{-10}$ & 数值一致 & 通过 \\
        \bottomrule
    \end{tabular}
\end{table}

表~\ref{tab:validation} 表明，能量平衡和储电量递推残差均远低于设定阈值，储能全过程没有越过安全边界；问题三加入互斥约束后，同时充、放电量仅处于数值舍入量级。与此同时，按“计划购电、调整交易、紧急购电”三部分重新计算全年费用，与模型目标值保持一致。由此可以确认，论文中的费用结论不是由约束遗漏或单位换算误差造成的，输出到五个结果工作簿的数据也与逐时求解明细一致。

\subsection{时间因果性的反事实扰动检验}

物理上可行的方案仍可能在信息使用上不现实。例如，如果 $0{:}00$ 制定的全天计划会随着“当天尚未发生的真实负荷”改变，那么即使费用更低，也不可能在实际运行中提前获得。因此，本文采用反事实扰动的方法检查非预见性：保持某个决策节点之前的全部历史信息不变，只人为改变该节点之后才会发生的数据；若该节点的计划没有使用未来信息，则其决策应保持不变。

具体检验包括五种情形：在问题二中分别改变当天尚未实现的真实负荷和真实光伏，检查 $0{:}00$ 计划是否保持不变；在问题三中改变 $6{:}00$ 之后的真实源荷，检查 $6{:}00$ 方案是否只由当时已知信息决定；在问题四的因果价格扩展中改变未来实际电价，检查当前节点计划是否不受影响；最后仅启用 $0{:}00$ 和 $6{:}00$ 两个节点，检查 $12{:}00$ 以后是否继续执行 $6{:}00$ 已经生效的方案，而不是错误退回日初计划。

上述检验全部满足预期，计划变量的最大数值偏差低于 $10^{-10}$。特别地，同样的未来价格扰动会改变“附件4直接使用”情形的结果，却不会改变因果价格扩展中扰动发生之前的计划，这说明结果不变确实来自信息隔离，而不是扰动本身没有作用。该检验从时间顺序上验证了式~\eqref{eq:causality}：本文主策略在每个节点都只使用当时已经获得的信息。

\subsection{调整结算与滚动计划的一致性检验}

问题三的费用规则会直接影响最优调整方向，因此结算公式不能只依赖程序内部计算，还需要用简单算例进行独立核对。取正常电价 $\pi=1$ 元/kWh，当日初计划为 $900$ kWh、执行前调整为 $700$ kWh 时，叠加结算式~\eqref{eq:settle-add} 得到 $900+0.5\times200=1000$ 元，差额结算式~\eqref{eq:settle-rep} 得到 $700+0.5\times200=800$ 元；当计划由 $700$ kWh 上调至 $900$ kWh 时，两种解释均得到 $700+1.5\times200=1000$ 元。手工计算与程序逐项回代一致。

更重要的是，两种结算方式改变的不只是“最后怎么算钱”，还会改变模型愿意向上或向下调整的程度。因此本文对两种口径分别从头求解，而不是把同一条调度曲线套入两套价格公式。随机生成的源荷与价格序列进一步表明，各自重新优化后的费用均不高于将另一口径最优方案直接代入本口径得到的费用，说明结算规则已经真正进入优化决策。

滚动执行方面，同一未来时段可能先后被 $0{:}00$、$6{:}00$、$12{:}00$ 和 $18{:}00$ 多次规划，但实际只执行最后一次生效的方案。对关闭部分节点的方案进行逐段追踪后，未发现“节点关闭后退回日初计划”的现象；各执行块均从上一阶段的实测储电量出发，并延续最近一次有效计划。这保证了问题三的节点消融比较只改变新增信息，而不同时改变计划延续规则。

\subsection{终端储电量策略的敏感性}
\label{sec:terminal}

问题一明确要求 $0{:}00$ 与 $24{:}00$ 储电量相同，而问题二、三没有给出同样的硬约束。若完全取消日末约束，单日优化可能把电池在窗口末端尽量放空，从而把本应由下一日承担的能量成本转移到当前日，形成有限时域的“末端透支”。因此主模型采用日循环策略，即计划层要求日末储电量回到当日实际初值；同时设置无终端约束的对照组，考察这一假设对结果的影响。

两种终端策略的费用与储电量变化如图~\ref{fig:sensitivity} 所示。日循环策略全年费用为 \SensTerminalCycle\ 元；取消终端约束后费用降为 \SensTerminalFree\ 元，变化幅度为 \SensTerminalPct\%，但全年最终储电量降至 \SensTerminalFreeSoc\ kWh，日末储电量均值也仅为 \SensTerminalFreeMean\ kWh，已经接近 $1200$ kWh 的安全下限。

\begin{figure}[htbp]
    \centering
    \includegraphics[width=0.98\textwidth]{figures/sensitivity.pdf}
    \caption{终端储电量策略与储能效率口径的敏感性}
    \label{fig:sensitivity}
\end{figure}

图~\ref{fig:sensitivity} 说明，无终端约束的费用优势主要来自“把电池留空”，而不是调度效率真正提高。虽然其日末 SOC 标准差为 \SensTerminalFreeStd\ kWh，小于日循环策略的 \SensTerminalCycleStd\ kWh，但这是因为大量日末状态集中在下限附近，不能据此认为运行更稳定。考虑到全年连续运行需要为下一日保留可比的储能条件，本文保留日循环策略作为主口径。

\subsection{储能效率口径的敏感性}
\label{sec:efficiency}

题目给出“充放电效率为 $90\%$”，存在两种常见理解：一是充电效率和放电效率分别为 $0.9$，对应往返效率 $0.81$；二是把 $90\%$ 理解为整个充放电过程的往返效率，此时两侧效率可取 $\sqrt{0.9}\approx0.9487$。为判断核心结论是否依赖这一文字解释，在其他条件不变的情况下分别重算，结果见表~\ref{tab:efficiency}。

\begin{table}[htbp]
    \caption{储能效率口径敏感性}
    \label{tab:efficiency}
    \centering
    \zihao{-5}
    \begin{tabular}{l r r r}
        \toprule
        效率口径 & 全年总费用/元 & 储能吞吐量/kWh & 紧急购电量/kWh \\
        \midrule
        口径A：$\eta_c=\eta_d=0.90$（往返 $0.81$） & \SensEffEOneCost & \SensEffEOneThroughput & \SensEffEOneEm \\
        口径B：$\eta_c=\eta_d=\sqrt{0.9}$（往返 $0.90$） & \SensEffETwoCost & \SensEffETwoThroughput & \SensEffETwoEm \\
        \bottomrule
    \end{tabular}
\end{table}

由表~\ref{tab:efficiency} 可见，提高往返效率会降低全年费用，这是因为单位充入电量能够在以后释放更多有效电能，储能在峰谷转移和偏差缓冲中的作用随之增强。但两种解释下，各问的策略方向和主要比较结论均保持一致，因此“储能能够降低高价购电、滚动信息能够减少偏差成本”的结论对效率口径具有稳定性。本文采用较保守的 $\eta_c=\eta_d=0.9$ 作为主模型参数。

\subsection{误差来源与适用边界}

模型剩余误差主要来自三个方面。第一，问题二的负荷和光伏只能依据历史序列预测，突发天气或异常负荷会使点预测和风险边际滞后；第二，附件3的小时级预报需要插值到 $10$ min，插值只能恢复时间分辨率，不能创造原预报中不存在的短时波动信息；第三，问题四正式结果按题意直接使用附件4电价，而真实市场中未来完整价格曲线未必能够提前获得。本文分别通过历史误差滚动更新、实时储能反馈、预报节点消融以及因果价格扩展对这些影响进行量化。它们不会破坏本文数值结果的内部一致性，但决定了模型向真实长期运行场景推广时仍需进一步引入气象、负荷和价格等外生预测信息。
